\section{模型评价与推广：数值答案及适用范围}

硅全量厚度为$2.17\,\mu\mathrm{m}$，非概率条件包络为$1.72$--$18.16\,\mu\mathrm{m}$。上一章的重估比较显示，厚度分支会随光学参数和窗口设置移动，本章据此评价双角约束的作用及后续标定方向。
% src:求解/问题3/结果/回炉轮3候选/仲裁闭环/20260910_131501_81623_建模公式/厚度结果.json:核心指标.硅全量完整往返厚度_um;核心指标.硅折1完整往返条件厚度_um;核心指标.硅折2完整往返条件厚度_um;核心指标.硅已计算条件包络下限_um;核心指标.硅已计算条件包络上限_um;核心指标.碳化硅正式基准厚度_um;核心指标.碳化硅正式条件范围下限_um;核心指标.碳化硅正式条件范围上限_um

碳化硅厚度为$10.80\,\mu\mathrm{m}$，厚度重估范围为$0.64$--$18.11\,\mu\mathrm{m}$（图\ref{fig:20}）。碳化硅沿用两束基准，未施加多光束修正（表\ref{tab:08}）。
% src:求解/问题3/结果/回炉轮3候选/仲裁闭环/20260910_131501_81623_建模公式/厚度结果.json:核心指标.硅全量完整往返厚度_um;核心指标.硅折1完整往返条件厚度_um;核心指标.硅折2完整往返条件厚度_um;核心指标.碳化硅正式基准厚度_um;核心指标.碳化硅正式条件范围下限_um;核心指标.碳化硅正式条件范围上限_um
% src:求解/问题3/结果/条件判定.json:碳化硅条件分支.触发

区段重估使硅厚度落在相距较远的分支；碳化硅的双角留段误差较二次趋势降低。完整往返场的高阶收益跨折反向，说明增加返回光并未稳定改善谱形预测。表\ref{tab:09}保留两问反射率测试覆盖的分子与分母，曲线预测区间仍有覆盖缺口。
% src:求解/问题3/结果/回炉轮3候选/仲裁闭环/20260910_131501_81623_建模公式/厚度结果.json:核心指标.硅折1完整往返条件厚度_um;核心指标.硅折2完整往返条件厚度_um
% src:求解/问题2/结果/验证.json:正式主方法相对正式二次趋势改善_百分比
% src:求解/问题3/结果/回炉轮3候选/仲裁闭环/20260910_131501_81623_建模公式/厚度结果.json:分材料验证.硅.逐折汇总.0.平均标准化均方根误差.两束;分材料验证.硅.逐折汇总.0.平均标准化均方根误差.完整往返;分材料验证.硅.逐折汇总.0.平均标准化均方根误差.训练均值;分材料验证.硅.逐折汇总.0.平均标准化均方根误差.二次趋势;分材料验证.硅.逐折汇总.1.平均标准化均方根误差.两束;分材料验证.硅.逐折汇总.1.平均标准化均方根误差.完整往返;分材料验证.硅.逐折汇总.1.平均标准化均方根误差.训练均值;分材料验证.硅.逐折汇总.1.平均标准化均方根误差.二次趋势;分材料验证.碳化硅.逐折汇总.0.平均标准化均方根误差.两束;分材料验证.碳化硅.逐折汇总.0.平均标准化均方根误差.完整往返;分材料验证.碳化硅.逐折汇总.0.平均标准化均方根误差.训练均值;分材料验证.碳化硅.逐折汇总.0.平均标准化均方根误差.二次趋势;分材料验证.碳化硅.逐折汇总.1.平均标准化均方根误差.两束;分材料验证.碳化硅.逐折汇总.1.平均标准化均方根误差.完整往返;分材料验证.碳化硅.逐折汇总.1.平均标准化均方根误差.训练均值;分材料验证.碳化硅.逐折汇总.1.平均标准化均方根误差.二次趋势
% src:求解/问题2/结果/区间覆盖.json:测试经验覆盖率
% src:求解/问题3/结果/回炉轮3候选/仲裁闭环/20260910_131501_81623_建模公式/厚度结果.json:预测区间.4.经验覆盖率;预测区间.4.平均区间宽度_比例;预测区间.5.经验覆盖率;预测区间.5.平均区间宽度_比例;预测区间.6.经验覆盖率;预测区间.6.平均区间宽度_比例;预测区间.7.经验覆盖率;预测区间.7.平均区间宽度_比例;预测区间.20.经验覆盖率;预测区间.20.平均区间宽度_比例;预测区间.21.经验覆盖率;预测区间.21.平均区间宽度_比例;预测区间.22.经验覆盖率;预测区间.22.平均区间宽度_比例;预测区间.23.经验覆盖率;预测区间.23.平均区间宽度_比例;预测区间.36.经验覆盖率;预测区间.36.平均区间宽度_比例;预测区间.37.经验覆盖率;预测区间.37.平均区间宽度_比例;预测区间.38.经验覆盖率;预测区间.38.平均区间宽度_比例;预测区间.39.经验覆盖率;预测区间.39.平均区间宽度_比例;预测区间.52.经验覆盖率;预测区间.52.平均区间宽度_比例;预测区间.53.经验覆盖率;预测区间.53.平均区间宽度_比例;预测区间.54.经验覆盖率;预测区间.54.平均区间宽度_比例;预测区间.55.经验覆盖率;预测区间.55.平均区间宽度_比例;预测区间.4.名义覆盖率
% src:求解/问题3/结果/回炉轮3候选/仲裁闭环/20260910_131501_81623_建模公式/厚度结果.json:逐块评分.4.测试点数;逐块评分.5.测试点数;逐块评分.6.测试点数;逐块评分.7.测试点数;逐块评分.20.测试点数;逐块评分.21.测试点数;逐块评分.22.测试点数;逐块评分.23.测试点数;逐块评分.36.测试点数;逐块评分.37.测试点数;逐块评分.38.测试点数;逐块评分.39.测试点数;逐块评分.52.测试点数;逐块评分.53.测试点数;逐块评分.54.测试点数;逐块评分.55.测试点数

现有双角观测缺少折射率曲线、相干长度及仪器分辨率标定。下文分别讨论双角约束的优点、光学参数互补的缺点，以及新增角度、标准厚度样片与光学标定的推广方向。
% src:交接/数据档案.json:总览.每块晶圆入射角度;未随附件提供的建模输入
% src:交接/结果声明_问题3.json:置信.理由
